\documentclass[../main.tex]{subfiles}
\begin{document}
\setlist[itemize]{topsep=2pt,itemsep=1pt,parsep=0pt,partopsep=0pt}
\setlist[itemize,2]{topsep=1pt,itemsep=0pt,parsep=0pt,partopsep=0pt}
\setlength{\parskip}{3pt}

\chapter{Week 6 — Reporting and Writing}
\label{ch:week6}

\textit{Snapshot and part summaries: see the overview on page~\pageref{ov:week6}.}

\section{Part 1 — Technical Reports}
\label{sec:reports}

\begin{itemize}
  \item Communicating the work is part of the job, not an addition to it.
  \item A \keyterm{technical report} conveys specific information about a technical subject to a specific audience for a specific and practical purpose.
  \item \keyterm{Research Methodologies report}: conveys your research plan to a knowledgeable peer in order to organise research activities.
  \item \keyterm{Report structure} — head (introduction, proportional to the body), body (connected sections), feet (conclusions that add no new information).
  \item RM report body: research questions (\autoref{sec:objectives-rqs}), methods with tools and expected results, planning (\autoref{sec:project-mgmt}).
  \item Write little and often (snack writing); alternate thinking and writing rather than separating them.
\end{itemize}

\section{Part 2 — Referencing}
\label{sec:referencing}

\begin{itemize}
  \item Use at least three document types including at least one book; at most two websites (or 10\% of the list).
  \item Most sources: peer-reviewed journal articles, mostly under ten years old (\autoref{sec:research-ideas}).
  \item Choose one reference style and keep to it throughout.
  \item \keyterm{Author-date (APA)}: in-text as \texttt{(Author et al., 2020)} or \texttt{Author (2020)}; list ordered alphabetically. Advantage: in-text identifies the source. Disadvantage: interrupts prose.
  \item \keyterm{Numbered (AIAA)}: in-text as \texttt{[1]}; list ordered by appearance. Advantage: clean text. Disadvantage: in-text marker carries no information.
  \item Reference checklist: every borrowed illustration and statement is cited; quotations in quotation marks; paraphrases in own words; list and text agree exactly in both directions.
\end{itemize}

\section{Part 3 — Writing Style}
\label{sec:writing-style}

\begin{itemize}
  \item Define the target reader first — what they want and what they need the information for.
  \item Readers scan: titles $\to$ figures $\to$ conclusions $\to$ body. Those elements must be self-contained.
  \item Layout: clean, well-placed high-resolution images; block diagrams often convey more than detailed explanations.
  \item \keyterm{4Cs}: Clear (terms defined, understood by audience), Comprehensive (all necessary information present), Concise (no repetition or padding), Correct (factually and grammatically sound).
  \item Graph conventions: tick marks with major gridlines at start and end; non-overlapping legend for multiple lines; markers same colour as line; straight or short-dashed lines; consistent units and fonts.
  \item \keyterm{AI writing tools}: use responsibly — generation can be redundant, brainstorming is backward-looking, generated references are error-prone.
  \item Final checklist: spelling and grammar; one or two ideas per paragraph; consistent format and nomenclature; acronyms defined on first use; every figure readable in PDF, captioned, and clear without the text; title and conclusions stand alone.
\end{itemize}

\end{document}
